%
%      Brno University of Technology
%      Faculty of Information Technology
%
%      MSc Thesis	2009/2010
%
%      Design of S-Boxes Using Genetic Algorithms
%
%
%      Bedrich Hovorka
%
%
%
\section{Definition of S-boxes}
A substitution box (S-box) is a component of a cipher whose output, as the name suggests,
is a substituted input. S-boxes bear most of the responsibility for nonlinearity
and randomness of the output of the entire cipher. The security of most block ciphers (especially those based on Feistel networks)
depends on the properties of S-boxes used in individual rounds. These will be described below.
Usually an S-box is depicted as a \uv{box} with individual inputs and outputs, where
individual bits are usually shown in order and numbering, as can be seen in Fig. \ref{sboxSchema}.

\insertimage[0.25\textwidth]{sboxSchema}{S-box schema}

Formally, an S-box can be defined as follows:
\begin{definition}\cite{practicalSboxDesign}
An $n \times m$ S-box $S$ is a mapping $S: \{0,1\}^n \rightarrow \{0,1\}^m$.
$S$ can be represented as $2^n$ $m$-bit numbers, denoted $r_0, \dots, r_{2^n-1}$,
in which case $S(x) = r_x$ , $0 \leq x < 2^n$ and $r_i$ are the rows
of the S-box. Alternatively, $S(x) = \left[ c_{m−1}(x) \, c_{m−2}(x) \dots c_{0}(x) \right] $
where $c_i$ is a fixed Boolean function $c_i : \{0,1\}^n \rightarrow \{0,1\} \, \forall i$;
these are the columns of the S-box.
\label{defSbox}
\end{definition}
The size of S-boxes is often written as $n \times m$.
Large S-boxes \cite{bigSboxes} are generally considered more secure
than small ones, but at the cost of greater implementation and validation complexity.
S-boxes with small input and large output have very good linear approximation.
For S-boxes with sufficiently larger output than input, at least one
perfect linear approximation is guaranteed.

\section{Security Criteria for S-boxes}\label{secCriterions}
Criteria that can be used as\footnote{or according to which can be constructed}
heuristic or evaluation functions can be used to search the state space.
Almost all are derived from (or have some relationship to)
linear and differential cryptanalysis.
They are derived from possible attacks on ciphers, some can be mathematically proven to provide resistance against
these attacks, while others are rather believed to work.

\subsection{SAC}\label{critSAC}
The Strict Avalanche Criterion \cite{yokohama} is based on the so-called \uv{avalanche effect}.
For S-boxes, we examine how many output bits change when one input bit changes. When testing
the change of one bit for all possible inputs, the average of changes at the output should be half the number of output
bits.

\begin{definition}[Avalanche effect] A function $f: \{0,1\}^n \rightarrow \{0,1\}^m$
exhibits the avalanche effect if and only if:
  \begin{equation}
    \sum\limits_{x \in \{0,1\}^n} wt(f(x) \oplus f(x \oplus c^{(n)}_i)) = m 2^{n-1}
  \end{equation}
for $\forall i: 1 \leq i \leq n$, where $wt(x)$ is the Hamming weight\footnote{number of one-bits}
and $c^{(n)}_i$ is an $n$-dimensional vector with $wt(c^{(n)}_i) = 1$ with position $i$ set.
\end{definition}

From the Avalanche effect, SAC is defined:

\begin{definition}[SAC, SAC of 0-th order] A function $f: \{0,1\}^n \rightarrow \{0,1\}^m$ satisfies SAC if
for $\forall i: 1 \leq i \leq n$ the following equations hold:
  \begin{equation}
    \sum\limits_{x \in \{0,1\}^n} f(x) \oplus f(x \oplus c^{(n)}_i) = (2^{n-1}, 2^{n-1}, \dots, 2^{n-1})
  \end{equation}
  Every function $f$ satisfying SAC is called a \uv{strong S-box}
\end{definition}

If a function satisfies SAC, each of its output bits should change with half probability,
every time one input bit is changed. If any output depends on only a few input
bits, then an attacker can use these relationships to find the key.
SAC has been further extended:

\begin{definition}[SAC of 1st order] A function $f: \{0,1\}^n \rightarrow \{0,1\}^m$ satisfies SAC of 1st order
if and only if:
  \begin{itemize}
    \item $f$ satisfies SAC of 0-th order
    \item every function obtained from $f$ by keeping one $i$-th bit constant and equal to $c$
      satisfies SAC for every $i \in \{1,2,\dots,n\}$ and $c = 0$ or $c = 1$.
  \end{itemize}
\end{definition}

In general, the extension can be defined recursively, up to order $k \leq n - 2$:

\begin{definition}[SAC of $k$-th order] A function $f: \{0,1\}^n \rightarrow \{0,1\}^m$ satisfies
SAC of $k$-th order if and only if:
  \begin{itemize}
    \item $f$ satisfies SAC of $(k - 1)$-th order
    \item every function obtained from $f$ by keeping $k$ of its bits constant
    (for all possible combinations of positions and values) satisfies SAC
  \end{itemize}
\end{definition}

\begin{table}
\begin{center}
  \begin{tabular}{| c | c | c | c | c |}
  \hline
  Function number & $00$ & $01$ & $10$ &  $11$\\ \hline
             1 &  $0$ &  $0$ &  $0$ &   $1$\\ \hline
             2 &  $0$ &  $0$ &  $1$ &   $0$\\ \hline
             3 &  $0$ &  $1$ &  $0$ &   $0$\\ \hline
             4 &  $1$ &  $0$ &  $0$ &   $0$\\ \hline
             5 &  $0$ &  $1$ &  $1$ &   $1$\\ \hline
             6 &  $1$ &  $0$ &  $1$ &   $1$\\ \hline
             7 &  $1$ &  $1$ &  $0$ &   $1$\\ \hline
             8 &  $1$ &  $1$ &  $1$ &   $0$\\ \hline
  \end{tabular}
  \caption{Basic functions satisfying MOSAC \cite{yokohama}}
  \label{miniSAC}
\end{center}
\end{table}

SAC of the highest possible order $k = n - 2$ has a special designation MOSAC\footnote{from English \uv{maximal order}}.
In table \ref{miniSAC}, the rows list the basic Boolean functions satisfying MOSAC.
These functions have a 2-bit input listed in the columns and satisfy SAC of 0-th order. The columns represent all
input combinations. As can be seen, for each function the outputs are the same for three bits and one
differs.

\begin{theorem}
A strong S-box is neither linear nor affine.
\end{theorem}
\begin{proof}
The proof is given in \cite{yokohama} on page 25
\end{proof}

\subsection{Bent functions}
In \cite{wikiBentFunction} combinatorics, a \uv{bent function} is a special type of Boolean function.
\uv{Bent functions} are so named because they are different
from linear and affine functions.
Therefore they have been extensively studied primarily for their application in cryptography,
although their spectrum of applications is broad also in coding theory and combinational circuit design.
The definition can be extended in many ways.

\begin{definition}\cite{yokohama} A Boolean function $g(w): \{0,1\}^n \rightarrow \{0,1\}$,
$n = 2l, l \in \mathcal{N}$ is a bent function,
if every coefficient
of the Fourier transform\footnote{this special case is more often called the Walsh transform}
$G(w)$ of the Walsh function $(-1)^{g(w)}$ defined
for all $w \in \{0,1\}^n$
\begin{equation}
  G(w) = \frac{1}{\sqrt{2^n}}\sum\limits_{x \in \{0,1\}^n}(-1)^{g(x)+x \cdot w}
\end{equation}
has unit magnitude
\begin{equation}
  |G(w)| = 1
\end{equation}
\end{definition}

\begin{theorem}[Relationship between bent functions and functions satisfying SAC]
\label{bentTheorem}
Let $\mathcal{A}_n $ denote the set of all $n$-bit Boolean functions,
$\mathcal{B}_n $ the set of $n$-bit bent functions
and $\mathcal{S}_n $ the set of $n$-bit Boolean functions satisfying SAC.
And in particular we can denote the set of $n$-bit Boolean functions satisfying
MOSAC as $\mathcal{S}^{\max}_n$. The relationship between these sets for even $n$ can be
expressed as
\begin{equation}
  \mathcal{S}^{\max}_n \subseteq \mathcal{B}_n \subseteq \mathcal{S}_n \subset \mathcal{A}_n
\end{equation}
\end{theorem}

\begin{proof}
The proof is given in \cite{yokohama} on page 31
\end{proof}
This can be used to narrow the search space when looking for functions satisfying SAC.
For S-boxes with multiple outputs, from definition \ref{defSbox} we can assume that each output is
a Boolean function and construct an evaluation function accordingly. For illustration, table
\ref{bentSetTable} shows the sizes of some sets.

\begin{table}
\begin{center}
  \begin{tabular}{| c | c | c | c | c | c |}
  \hline
  n & 2 & 3 & 4 & 5 & 6 \\ \hline
  $|\mathcal{A}_n|$ & $16$ & $256$ & $65536$ & $2^{32}$ & $2^{64}$ \\ \hline
  $|\mathcal{S}_n|$ & $8$ & $64$ & $4128$ & ? & ? \\ \hline
  $|\mathcal{B}_n|$ & $8$ & -- & $896$ & -- & $2^{32.3}$ \\ \hline
  $|\mathcal{S}^{\max}_n|$ & $8$ & $16$ & $32$ & $64$ & $128$ \\ \hline
  \end{tabular}
  \caption{Size of sets from theorem \ref{bentTheorem} \cite{yokohama}}
  \label{bentSetTable}
\end{center}
\end{table}

\subsection{Maximum deviation from linear expressions}
To calculate the nonlinearity of a Boolean function $f$ we can use the formula from \cite{yokohama},
based on the Hamming distance from the set of affine functions and on the Walsh transform
$\hat{F}$ of the function $\hat{f}: \{0,1\}^n \rightarrow \{1, -1\}$, similarly as for SAC:
\begin{eqnarray}
  \hat{f}(x) & = & (-1)^{f(x)} \\
  \hat{F}(w) & = & \sum\limits_{x \in \{0,1\}} f(x) \cdot (-1)^{x \cdot w}\\
  \delta(f) & = & 2^{n-1} - \frac{1}{2} \max_{w \in \{0,1\}}|\hat{F}(w)|
\end{eqnarray}

But for multiple outputs it will be simpler to calculate the maximum deviation of S-box $q$
according to the formula from \cite{twofish}:
\begin{equation}
LP_{\max}(q) = \max_{a, b \neq 0} \left( 2 Pr_X[X \cdot a = q(X) \cdot b] - 1 \right)^2
\end{equation}

This formula is based on the procedure of first
creating a complete frequency table of all linear approximations corresponding to the form of equation \ref{eqLinearAprox}.
Rows resp. columns correspond to individual parameters $a$ resp. $b$,
$\forall a, b$ where $0 \leq a < n$ and $0 \leq b < m$\footnote{$n$ and $m$ are according to def. \ref{defSbox}},
Each cell of the table \cite{tutorialLADA} represents
the probability of occurrence (or number of matches) with the linear expression corresponding to parameters $a$ and $b$. And from them
the largest linearity is selected and then used when searching for an S-box that has
this maximum linearity as small as possible.
The deviation is normalized in the formula to the interval $\langle 0,1 \rangle$.
As follows from the text above, its magnitude is of interest, not the sign.

\pagebreak

\subsection{Maximum probability from differentials}
To calculate the maximum probability of the differential of S-box $q$
we use this formula \cite{twofish}, derived from differential cryptanalysis:
\begin{equation}
  DP_{\max}(q) = \max_{a \neq 0 , b} Pr_X[q(X \oplus a) \oplus q(X) = b]
\end{equation}
This formula is again based on a table \cite{tutorialLADA} of occurrence counts
(probability distribution) of differences, in which each row represents values
$\Delta X = a$ and columns represent values $\Delta Y = b$.
Each cell of the table represents the number of cases when the
output difference equals $b$, if the input difference is $a$. Note that for example the case excluded
from calculation having parameters ($a = 0$, $b = 0$), always has the highest probability in the table,
which however does not belong to a real difference. Similarly as for the criterion $LP_{\max}$, the maximum
probability is used from the table for minimization.
In \cite{practicalSboxDesign} this criterion is called the \uv{XOR table}.

\subsection{Branching factor}\label{secBF}
The encryption process should diffuse the encoded information as much as possible,
so that the search space is as complex as possible for the attacker. Imagine the
search space in the form of a tree. The minimum number of successors of each node
determines the lower bound of search complexity, which depends on it exponentially.

The branching factor (Branching Factor, Branch Number) is used in \cite{rijndael} as the main
criterion in the linear step of the cipher named \uv{MixColumn} by bytes. However, nothing prevents
testing it also on a nonlinear element by bits.
Let $f$ be an S-box according to \ref{defSbox}, $wt(x)$ denotes the Hamming weight.
The branching factor is a measure of diffusion strength:

\begin{definition}
  The branching factor of function $f$ is
  \begin{eqnarray}
    \min_{\Delta x \neq 0, x} \left( wt(\Delta x) + wt(\Delta f(x)) \right)\\
    \Delta f(x) = f(x) \oplus f(x \oplus \Delta x)
  \end{eqnarray}
\end{definition}
At the input and output are in fact differences, the minimum sum of the number of changes at the input and output together is calculated.
This criterion is used for maximization.

\subsection{Degree of algebraic polynomial}
The function that describes an S-box can be approximated not only by linear expressions in $GF(2)$ arithmetic,
but also interpolated by polynomials in this arithmetic for each output bit. The polynomial for one of the bits
of output $y_i$ can be defined using the formula:
\begin{equation}
  y_i(x) = \sum\limits_{a \in \mathcal{A}} \left(\prod\limits_{j=0}^{n} a_j \cdot z_j \right)
\end{equation}
Where $z_0$ represents 1, $z_j = x_{j - 1}$ are individual input bits\footnote{in the previous text $x$ is counted from 0},
$\mathcal{A}$ is the set of parameters determining (turning on individual bits of) the polynomial.
The degree of the polynomial for one output bit is the maximum Hamming weight from the set $\mathcal{A}$ (the bit for one is not included).
For the entire S-box, the minimum degree from the polynomials for outputs can then be used as a criterion:
\begin{equation}
  \min\limits_{0 \leq i < m} \max\limits_{a \in \mathcal{A}_i} \left( wt(a_{[1 \cdots n]}) \right)
\end{equation}

\pagebreak
A $4 \times 4$ S-box given by the sequence
\begin{verbatim}
  [ 7 11 15 2 8 4 1 3 6 9 14 5 12 0 13 10 ]
\end{verbatim}
\nopagebreak
is expressed by these third-degree polynomials:
\nopagebreak
\begin{eqnarray}
  y_0 & = & 1 + a*b + c + b*c + a*b*c + d + a*d + a*b*d + c*d + a*c*d\\
  y_1 & = & 1 + c + a*b*c + a*d + a*c*d\\
  y_2 & = & 1 + a + c + a*b*c + a*b*d + c*d\\
  y_3 & = & a + b + c + a*b*c + b*c*d
\end{eqnarray}
where $a = x_0$ is the first input bit, etc. The operations $+$ and $*$ are defined in mod-2 arithmetic.
And they express how much the output bits depend on the input bits. If a polynomial has the highest
possible degree, the output bit depends on all inputs. For the entire S-box, therefore, the minimum
degree from all polynomials expressing the outputs is chosen. Another possibility
could be the number of variables used, but addition provides information about linearity.

\subsection{Bijectivity}
Some cryptosystems\footnote{if $n=m$ then in that case logically perhaps all should} require
that the function represented by the S-box be bijective -- i.e., each input is mapped to
a unique output.
This can of course be simply achieved by encoding the problem as a permutation (see below), but if
we want to use a different type of representation, we must for example establish bijectivity as an additional criterion.
In \cite{yokohama} a constraint is further mentioned on constructing a bijective S-box
from Boolean functions satisfying MOSAC. Let there be a bijective function
$f: \{0,1\}^n \rightarrow \{0,1\}^n$, every linear combination of Boolean functions
from which it is composed must have a Hamming weight equal to half of the possible output combinations:
\begin{equation}
  wt(\sum\limits_{i=1}^{n} a_i f_i) = 2^{n-1}
\end{equation}
In other words: all linear deviations, according to expression \ref{eqLinearAprox} not containing
input bits, must be zero.

\subsection{Other criteria}
Other criteria that were not further used include for example the invertibility of function
$f$. Or for example the BIC criterion (bit independence) from \cite{practicalSboxDesign} or
\uv{Cross Correlation of Avalanche Variables} from \cite{yokohama}.
Both are derived from the avalanche effect.


\begin{table}
  \begin{longtable}{| l | p{0.1875\textwidth} | p{0.125\textwidth} | p{0.125\textwidth} | p{0.0625\textwidth} |}
  \hline
  Criterion & modification & minimum value & maximum value & optimization type \\ \hline
bent score & number of $w$ whose $|G(w)| = 1$ & $0$ & $m \cdot 2^n$ & $\max$ \\ \hline
SAC degree & ~ & -1 & $n - 2$ & $\max$ \\ \hline
differential probability & ~ & * $\frac{1}{2^n}$ & $1$ & $\min$ \\ \hline
linear bias probability & ~ & * $0$ & 1 & $\min$ \\ \hline
algebraic polynomial degree & ~ & $0$ & $n$ & $\max$ \\ \hline
%řád Lagrangeova polynomu & ~ & $0$ & $m$ & $\max$ \\ \hline
branch factor & ~ & 0 & $m+1$ &  $\max$ \\ \hline
 bijective score & similar to bent score & 0 & $m \cdot 2^n$ & $\max$ \\ \hline
  \end{longtable}
  \caption{Comparison of individual security criteria}
  \label{criterionComparisonTable}
\end{table}

\bigskip
In table \ref{criterionComparisonTable}, the individual criteria used are compared. The boundary
values of individual criteria are listed here. Some of them depend on the number of inputs and outputs.
For two of them, the minimum value is sought.
Asterisks mark theoretical values that
cannot be practically achieved.

\pagebreak % bude treba?
\section{Examples of further use of S-boxes}\label{secSboxUsing}
S-boxes are described for many ciphers.
Examples of the most well-known ones follow:
\subsection{S-box in DES}
In DES \cite{wikiSbox}, $6 \times 4$ S-boxes are used.
A good example is the S-box $S_5$ shown in table \ref{desSbox}:
\begin{table}[pt]
\begin{center}
  \begin{tabular}{| l | c | c | c | c |  c | c | c | c |}
  \hline
Outer bits	& \multicolumn{8}{|c|}{Inner 4 input bits} \\ \hline
~  & \bf 0000 & \bf 0001 & \bf 0010 & \bf 0011 & \bf 0100 & \bf 0101 & \bf 0110 & \bf 0111 \\ \hline
00 & 0010 & 1100 & 0100 & 0001 & 0111 & 1010 & 1011 & 0110 \\ \hline
01 & 1110 & 1011 & 0010 & 1100 & 0100 & 0111 & 1101 & 0001 \\ \hline
10 & 0100 & 0010 & 0001 & 1011 & 1010 & 1101 & 0111 & 1000 \\ \hline
11 & 1011 & 1000 & 1100 & 0111 & 0001 & 1110 & 0010 & 1101 \\ \hline
  \hline
~  & \bf 1000 & \bf 1001 & \bf 1010 & \bf 1011 & \bf 1100 & \bf 1101 & \bf 1110 & \bf 1111 \\ \hline
00 & 1000 & 0101 & 0011 & 1111 & 1101 & 0000 & 1110 & 1001 \\ \hline
01 & 0101 & 0000 & 1111 & 1010 & 0011 & 1001 & 1000 & 0110 \\ \hline
10 & 1111 & 1001 & 1100 & 0101 & 0110 & 0011 & 0000 & 1110 \\ \hline
11 & 0110 & 1111 & 0000 & 1001 & 1010 & 0100 & 0101 & 0011 \\ \hline
  \end{tabular}
  \caption{DES S-box $S_5$}
  \label{desSbox}
\end{center}
\end{table}
For a 6-bit input, the 4-bit output is determined by selecting a row using the two outer bits (first and last),
the column is determined by the inner four bits. E.g. input \uv{011011} has outer bits \uv{01} and inner \uv{1101}.
The corresponding output is \uv{1001}.
The eight S-boxes of the DES algorithm were many years ago the subject of intensive study with suspicion of a back
door, a vulnerability known only to its creators, that could be placed in the cipher.
The design criteria for the S-boxes were finally published (Don Coppersmith, 1994),
after the rediscovery of differential cryptanalysis,
having been carefully improved in resistance to this attack.
Later research showed that precisely these small changes could have significantly weakened DES.

\subsection{S-box in AES}
The Rijndael cipher\cite{rijndael, wikiRijnDaelSbox} uses S-boxes in a phase called \uv{ByteSub}.
Here we are dealing with $8 \times 8$ S-boxes,
i.e. byte to byte. It was designed with regard to resistance to linear and differential cryptanalysis,
to have the highest polynomial degree be invertible and easily describable.
Therefore they chose the following affine transformation, modular multiplication
followed by addition:
\begin{equation}
b(x) = (x^7 + x^6 + x^2 + x) + a(x)(x^7 + x^6 + x^5 + x^4 + 1) \mod x^8 + 1
\end{equation}
Where $(b(a) \neq a)$ and $(b(a) \neq \bar{a} )$.

\subsubsection{S-boxes in Twofish}
Here \cite{twofish} they use $8 \times 8$ to construct $8 \times 32$.
These small S-boxes are determined by random search, in which they looked for permutations
with $LP_{\max} \leq \frac{1}{16}$ and $DP_{\max} \leq \frac{10}{256}$.
